\section{Affine trick}

\subsection{Space-side affine trick}

Let $\Lambda_K^{(d)}$ be the spatial form from
\eqref{eq:CZ-space-operator} in input dimension $d$. The space-side trick
has two useful forms. The first simply freezes the last coordinate of every
input. The second resolves the zero-sum constraint in affine coordinates and
is better suited to estimates.

For the symmetric form, let
\[
 d\mu_\perp(\vec t) \eqd \delta_0\left(\sum_{j=1}^{d+1}t_j\right)  \prod_{j=1}^{d+1}\dd t_j
\]
denote Lebesgue measure on $\vec{1}^\perp$ with this normalization.

\begin{proposition}[Space-side affine trick: symmetric slicing]
\label{prop:symmetric-spatial-affine-trick}
For Schwartz functions $f_1,\ldots,f_{d+1}$, one has
\[
\Lambda_K^{(d)}(f_1,\ldots,f_{d+1})
\begin{aligned}[t]
&=\int_{\vec{1}^\perp} \mathcal{S}_{\vec t,K}^{(d-1)} (F_{1,t_1},\ldots,F_{d+1,t_{d+1}}) \dd \mc{L}^{d-1}(\vec t),
\\
&=  \int_0^\infty r^{d-1}\int_{\mathbb{S}^{d-1}\cap\vec{1}^\perp}  \mathcal{S}_{\omega,K_r}^{(d-1)} (F_{1,r\omega_1},\ldots,F_{d+1,r\omega_{d+1}}) \dd\sigma^{d-2}(\omega)\dd r.
\end{aligned}
\]
where \(  F_{i,t}(x)\eqd f(x,t)\), \((x,t)∈\R^{d-1}\times\R\), $\dd\sigma^{d-2}$ is the surface measure on
$\mathbb{S}^{d-1}\cap\vec{1}^\perp$ given by \(\dd\mc{L}^{d-1}(r\omega)=r^{d-1}\dd r\dd\sigma_\perp(\omega)\), and $K_r(s)\eqd K(rs)$.
\end{proposition}

\begin{proof}
Write $u_j=(x_j,t_j)\in\mathbb{R}^{d-1}\times\mathbb{R}$.
The zero-sum constraint in the definition of $\Lambda_K^{(d)}$ is the product of the two constraints $\sum_jx_j=0$ and $\sum_jt_j=0$.
Moreover,
\[
D^{(d)}(u_1,\ldots,u_{d+1})
=\det\begin{pmatrix}
x_1&\cdots&x_{d+1}\\
t_1&\cdots&t_{d+1}\\
1&\cdots&1
\end{pmatrix}
=\mathcal{D}_{\vec t}^{(d-1)}(x_1,\ldots,x_{d+1}).
\]
Substitution into the defining integral gives the first identity.
The second follows from polar coordinates on $\vec{1}^\perp$ and the identity $\mathcal{D}_{r\omega}^{(d-1)}=r\mathcal{D}_{\omega}^{(d-1)}$.
\end{proof}

\begin{corollary}[Reduction from sliced bounds]
\label{cor:uniform-sliced-bound-reduction}
Let $p_1,\ldots,p_{d+1}\in(1,\infty]$ and $q\in[1,\infty]$ satisfy \(\frac{1}{q}+\sum_{j=1}^{d+1}\frac{1}{p_j}=1\).
Suppose that
\[
\bigl|\mathcal{S}_{\omega,K_r}^{(d-1)} (g_1,\ldots,g_{d+1})\bigr| \le C_{\vec p}(r,\omega) \prod_{j=1}^{d+1}\|g_j\|_{L^{p_j}(\R^{d-1})},
\qquad r>0,\quad \omega\in\mathbb{S}^{d-1}\cap\vec{1}^{\perp},
\]
and set
\[
\begin{aligned}
\mathfrak {C}_{\vec p,q}
&\eqd\left\|r^{d-1}C_{\vec p}(r,\omega) \prod_{j=1}^{d+1}|\omega_j|^{-1/p_j}\right\|_{ L^1_\omega(\mathbb{S}^{d-1}\cap\vec{1}^{\perp};L^q_r(0,\infty))}
\\
&= \int_{\mathbb{S}^{d-1}\cap\vec{1}^{\perp}} \bigl\|r^{d-1}C_{\vec p}(r,\omega)\bigr\|_{L^q_r(0,\infty)} \prod_{j=1}^{d+1}|\omega_j|^{-1/p_j} \dd\sigma_\perp(\omega).
\end{aligned}
\]
If \(\mathfrak {C}_{\vec p,q} <\infty\) then
\[
\bigl|\Lambda_K^{(d)}(f_1,\ldots,f_{d+1})\bigr| \lesssim_{d,\vec p} \mathfrak {C}_{\vec p,q} \prod_{j=1}^{d+1}\|f_j\|_{L^{p_j}(\R^d)}.
\]
\end{corollary}

\begin{proof}
Put $A_j(t)\eqd\|f_{j,t}\|_{L^{p_j}(\R^{d-1})}$.
The polar identity and Hölder's inequality in $r$ give, for almost every $\omega$,
\[
\begin{aligned}
& \int_0^\infty r^{d-1}C_{\vec p}(r,\omega) \prod_{j=1}^{d+1}A_j(r\omega_j)\,\dd r
\le \bigl\|r^{d-1}C_{\vec p}(r,\omega)\bigr\|_{L^q_r(0,\infty)} \prod_{j=1}^{d+1}|\omega_j|^{-1/p_j} \|A_j\|_{L^{p_j}(\R)}.
\end{aligned}
\]
Integration in $\omega$, followed by Fubini and the identity $\|A_j\|_{L^{p_j}(\R)}=\|f_j\|_{L^{p_j}(\R^d)}$, proves the claim.
\end{proof}

\begin{corollary}[The plane bound from symmetric slicing]
\label{cor:symmetric-slicing-plane}
Let $K$ be a smooth Calderón--Zygmund kernel on $\R$ and let $p_1,p_2,p_3\in(2,\infty)$ and \( \frac1{p_1}+\frac1{p_2}+\frac1{p_3}=1 \).
Assume that, for every smooth Calderón--Zygmund kernel $L$ on $\R$ whose Calderón--Zygmund seminorms are bounded by $1$, one has
\[
\left|\pv\int_{\R^2} h_1(x+t)h_2(x+\sigma t)h_3(x)L(t)\,\dd t\dd x\right|
\lesssim_{\vec p} \prod_{j=1}^3\|h_j\|_{L^{p_j}(\R)},
\]
uniformly in $\sigma\in\R$. Then
\[
\bigl|\Lambda_K^{(2)}(f_1,f_2,f_3)\bigr| \lesssim_{K,\vec p} \prod_{j=1}^3\|f_j\|_{L^{p_j}(\R^2)}.
\]
\end{corollary}

\begin{proof}
Fix $\omega\in\mathbb{S}^1\cap\vec{1}^{\perp}$ and let
\[
\nu\eqd\frac1{\sqrt3} (\omega_2-\omega_3,\omega_3-\omega_1,\omega_1-\omega_2).
\]
Then $(\omega,\nu)$ is an orthonormal basis of $\vec{1}^{\perp}$. Writing
$(x_1,x_2,x_3)=s\omega+t\nu$, one has
\[
\mathcal D_\omega^{(1)}(x_1,x_2,x_3)=\sqrt3\,t
\]
up to an irrelevant sign. Consequently,
\[
\mathcal S_{\omega,K_r}^{(1)}(g_1,g_2,g_3) =\frac1{\sqrt3}\int_{\R^2}\prod_{j=1}^3g_j(s\omega_j+t\nu_j) K(r\sqrt3\,t)\,\dd s\dd t.
\]
Away from the measure-zero set where some $\omega_j=0$, put
\[
G_j(u)\eqd g_j(\omega_ju),   \qquad \beta_j\eqd\frac{\nu_j}{\omega_j}.
\]
After translating and rescaling the $t$-variable, the last display is a
bilinear Hilbert transform with one free slope parameter.
The rescaled kernels have Calderón--Zygmund seminorms $O_K(r^{-1})$, uniformly in $\omega$.
Hence Theorem~\ref{thm:uniform-bht} gives
\[
\bigl|\mathcal S_{\omega,K_r}^{(1)}(g_1,g_2,g_3)\bigr| \lesssim_{K,\vec p}r^{-1} \prod_{j=1}^3|\omega_j|^{-1/p_j}\|g_j\|_{L^{p_j}(\R)}.
\]
Apply Corollary~\ref{cor:uniform-sliced-bound-reduction} with $q=\infty$ and
\[
C_{\vec p}(r,\omega)\lesssim_{K,\vec p} r^{-1}\prod_{j=1}^3|\omega_j|^{-1/p_j}.
\]
Since $rC_{\vec p}(r,\omega)$ is independent of $r$, the mixed norm condition becomes
\[
\int_{\mathbb{S}^1\cap\vec{1}^{\perp}} \prod_{j=1}^3|\omega_j|^{-2/p_j}\dd\sigma_\perp(\omega)<\infty.
\]
Each coordinate $\omega_j$ has only simple zeros on this circle, so the last integral is finite precisely when $2/p_j<1$ for every $j$.
This is exactly the assumption $p_j>2$.
\end{proof}

\subsection{Multiplier-side affine trick}
For each $v\in\R$, the reduced kernel in the decomposition below is the
one-variable distribution
\[
  \mu_v\eqd\mu_{\omega,v},
  \qquad
  \mu_v(t)\eqd\FT{K}\left(\frac vt\right),\qquad t\ne0.
\]
It is independent of $\omega$; the subscript $\omega$ only records the reduced form
in which this kernel is used.  For a distributional $K$, this expression,
as well as the formula below, is understood by regularization.

\begin{proposition}[Multiplier-side affine trick: spherical coarea form]
\label{prop:multiplier-affine-trick}
Let $g_1,\ldots,g_{d+1}\in\Sch(\R^d)$.  For
$\omega\in\mathbb S^{d-1}\cap\vec1^\perp$ and $u\in\R$,
write
\[
  g_{j,u}(x)\eqd g_j(x,u),\qquad x\in\R^{d-1}.
\]
Then
\[
\mc{M}_K^{(d)}(g_1,\ldots,g_{d+1}) =\frac12\int_\R \int_{\mathbb S^{d-1}\cap\vec1^\perp}\mathcal A_{\omega,\mu_v}^{(d-1)} \bigl(g_{1,\omega_1v},\ldots,g_{d+1,\omega_{d+1}v}\bigr) \dd\sigma_\perp(\omega) \dd v.
\]
\end{proposition}

\begin{proof}
We use the coarea identity
\[
\int F(\xi)\prod_{j=1}^{d+1}\dd\xi_j
    =\frac12\int_{\mathbb S^{d-1}\cap\vec1^\perp}\int F(\xi)\,
|\mc{D}^{(d-1)}_\omega(\xi_1,\ldots,\xi_{d+1})|\,
\delta_0\left(\sum_{j=1}^{d+1}\omega_j\xi_j\right)
\prod_{j=1}^{d+1}\dd\xi_j\dd\sigma_\perp(\omega).
\]
Indeed, for a full-rank configuration $(\xi_1,\ldots,\xi_{d+1})$, the
constraint above has exactly two solutions $\omega\in\mathbb {S}^{d-2}\cap\vec1^\perp$.
The two points account for the factor $1/2$, while the coarea Jacobian is $|\mc{D}^{(d-1)}_\omega(\xi_1,\ldots,\xi_{d+1})|$.
The rank-deficient set has measure zero.

  Partial Fourier transformation in $x\in\R^{d-1}$ gives
  \[
    \FT{g_{j,\omega_jv}}(\xi_j)
    =\int_\R\FT{g_j}(\xi_j,\tau_j)
    e^{2\pi i \omega_jv\tau_j}\dd\tau_j.
  \]
  On the support of the delta constraint in \((*)\), elementary row
  operations give
  \[
    D^{(d)}(\zeta_1,\ldots,\zeta_{d+1})
    =\mc{D}^{(d-1)}_\omega(\xi_1,\ldots,\xi_{d+1})
    \sum_{j=1}^{d+1}\omega_j\tau_j.
  \]
  The one-dimensional scaling identity
  \[
    \int_\R\widehat K\left(\frac vt\right)e^{2\pi ivs}\dd v
    =|t|K(ts),\qquad t\ne0,
  \]
  with $t=\mc{D}^{(d-1)}_\omega(\xi_1,\ldots,\xi_{d+1})$ and
  $s=\sum_j\omega_j\tau_j$, therefore yields
  \[
    |\mc{D}^{(d-1)}_\omega(\xi_1,\ldots,\xi_{d+1})|K\left(\mc{D}^{(d-1)}_\omega(\xi_1,\ldots,\xi_{d+1})
    \sum_{j=1}^{d+1}\omega_j\tau_j\right)
    =|\mc{D}^{(d-1)}_\omega(\xi_1,\ldots,\xi_{d+1})|
    K\bigl(D^{(d)}(\zeta_1,\ldots,\zeta_{d+1})\bigr).
  \]
  Applying the coarea formula to the remaining $\xi$-integral proves the formula.
  As usual, distributional kernels are obtained by regularization and a
  passage to the limit.
\end{proof}

The formula makes the change in homogeneity transparent.  If $K$ is
homogeneous of degree $-1$, then $\widehat K$ is homogeneous of degree $0$;
thus, for fixed $v$,
\[
  \mu_v(\lambda t)=\mu_v(t),\qquad \lambda>0.
\]
The $(-1)$-homogeneous multiplier is therefore resolved into reduced
determinantal forms with a $0$-homogeneous kernel.
\begin{corollary}[Reduction from multiplier slices]
\label{cor:uniform-multiplier-sliced-bound}
Let $p_1,\ldots,p_{d+1}\in(1,\infty]$ and $q\in[1,\infty]$ satisfy
\[
  \frac{1}{q}+\sum_{j=1}^{d+1}\frac1{p_j}=1.
\]
Suppose that
\[
  \bigl|\mathcal A_{\omega,\mu_v}^{(d-1)}
  (h_1,\ldots,h_{d+1})\bigr|
  \le C_{\vec p}^{\mathrm{mult}}(v,\omega)
  \prod_{j=1}^{d+1}\|h_j\|_{L^{p_j}(\R^{d-1})}
\]
for every $v\in\R$ and $\omega\in\mathbb S^{d-1}\cap\vec1^\perp$.  Set
\[
  \mathfrak C_{\vec p,q}^{\mathrm{mult}}
  \eqd\int_{\mathbb S^{d-1}\cap\vec1^\perp}
  \bigl\|C_{\vec p}^{\mathrm{mult}}(\cdot,\omega)\bigr\|_{L^q(\R)}
  \prod_{j=1}^{d+1}|\omega_j|^{-1/p_j}\,
  \dd\sigma_\perp(\omega).
\]
If $\mathfrak C_{\vec p,q}^{\mathrm{mult}}<\infty$, then
\[
  \bigl|\mc M_K^{(d)}(g_1,\ldots,g_{d+1})\bigr|
  \lesssim_{d,\vec p}
  \mathfrak C_{\vec p,q}^{\mathrm{mult}}
  \prod_{j=1}^{d+1}\|g_j\|_{L^{p_j}(\R^d)}.
\]
\end{corollary}

\begin{proof}
Put $A_j(u)\eqd\|g_{j,u}\|_{L^{p_j}(\R^{d-1})}$.  The formula in
Proposition~\ref{prop:multiplier-affine-trick}, the assumed bound, and
H\"older's inequality in $v$ give
\[
  \bigl|\mc M_K^{(d)}(g_1,\ldots,g_{d+1})\bigr|
  \lesssim_{d,\vec p}
  \mathfrak C_{\vec p,q}^{\mathrm{mult}}
  \prod_{j=1}^{d+1}\|g_j\|_{L^{p_j}(\R^d)}.
\]
Here we used $\|A_j\|_{L^{p_j}(\R)}=\|g_j\|_{L^{p_j}(\R^d)}$.
\end{proof}

\subsection{Projectively normalized slices and weighted norms}

There is a useful way to apply the two slicing reductions simultaneously to
all the inputs.  For a function \(g\) on
\(\R^{d-1}\times\R\), define its projectively normalized slices by
\[
  G_t(x)\eqd g(tx,t),\qquad t\ne0.
\]
Thus \(g_t(x)=g(x,t)=G_t(x/t)\).  Consequently, for
\(p\in(1,\infty)\),
\[
  \|g_t\|_{L^p(\R^{d-1})}
  =|t|^{(d-1)/p}\|G_t\|_{L^p(\R^{d-1})},
\]
and hence
\[
  \boxed{
  \|g\|_{L^p(\R^d)}^p
  =\int_\R\int_{\R^{d-1}}|G_t(x)|^p\,
    |t|^{d-1}\dd x\dd t.}
  \eqnum \label{eq:projective-weighted-norm}
\]
For \(p=\infty\), the corresponding statement is
\(\|g\|_\infty=\operatorname*{ess\,sup}_t\|G_t\|_\infty\).  Thus the ordinary
\(L^p\)-norm of \(g(x,t)\) is precisely the weighted \(L^p\)-norm of the
family \(G_t(x)\), with projective measure
\(|t|^{d-1}\dd x\dd t\).

Apply this construction to every input \(g_j\), and write
\[
  G_{j,t}(x)\eqd g_j(tx,t).
\]
We now absorb the inverse dilations \(x\mapsto x/(\omega_jv)\) and
\(x\mapsto x/(r\omega_j)\) into the operators themselves.  Define
\[
\begin{aligned}
  \widetilde{\mathcal A}_{\omega,v}^{(d-1)}
  (H_1,\ldots,H_{d+1})
  &\eqd
  \mathcal A_{\omega,\mu_v}^{(d-1)}
  \bigl(
    H_1(\cdot/(\omega_1v)),\ldots,
    H_{d+1}(\cdot/(\omega_{d+1}v))
  \bigr),\\
  \widetilde{\mathcal S}_{\omega,r}^{(d-1)}
  (H_1,\ldots,H_{d+1})
  &\eqd
  \mathcal S_{\omega,K_r}^{(d-1)}
  \bigl(
    H_1(\cdot/(r\omega_1)),\ldots,
    H_{d+1}(\cdot/(r\omega_{d+1}))
  \bigr).
\end{aligned}
\eqnum \label{eq:modified-affine-operators}
\]
These are the modified affine operators.  They act directly on fibers:
\[
  \mc M_K^{(d)}(g_1,\ldots,g_{d+1})
  =\frac12\int_\R\int_{\mathbb S^{d-1}\cap\vec1^\perp}
  \widetilde{\mathcal A}_{\omega,v}^{(d-1)}
  \bigl(
    G_{1,\omega_1v},\ldots,G_{d+1,\omega_{d+1}v}
  \bigr)
  \dd\sigma_\perp(\omega)\dd v,
\]
and
\[
  \Lambda_K^{(d)}(g_1,\ldots,g_{d+1})
  =\int_0^\infty r^{d-1}
  \int_{\mathbb S^{d-1}\cap\vec1^\perp}
  \widetilde{\mathcal S}_{\omega,r}^{(d-1)}
  \bigl(
    G_{1,r\omega_1},\ldots,G_{d+1,r\omega_{d+1}}
  \bigr)
  \dd\sigma_\perp(\omega)\dd r.
\]
Thus the arguments of the modified operators are the fibers themselves,
with no residual dilation in their arguments.

We next make explicit how absorbing the dilations changes the geometry.
For the multiplier-side operator, set
\[
  \mathcal R_\omega^{(d-1)}(\eta_1,\ldots,\eta_{d+1})
  \eqd
  \det\begin{pmatrix}
    \eta_1/\omega_1&\cdots&\eta_{d+1}/\omega_{d+1}\\
    \omega_1&\cdots&\omega_{d+1}\\
    1&\cdots&1
  \end{pmatrix}.
\]
The change of frequency variables
\(\eta_j=\omega_jv\xi_j\) in the definition of
\(\mathcal A_{\omega,\mu_v}^{(d-1)}\) gives
\[
\begin{aligned}
  \widetilde{\mathcal A}_{\omega,v}^{(d-1)}
  (H_1,\ldots,H_{d+1})
  ={}&|v|^{d-1}
  \int_{(\R^{d-1})^{d+1}}
  \prod_{j=1}^{d+1}\widehat H_j(\eta_j)\\
  &\quad\times
  \mu_v\left(
    v^{-(d-1)}\mathcal R_\omega^{(d-1)}
    (\eta_1,\ldots,\eta_{d+1})
  \right)\\
  &\quad\times
  \delta_0\left(\sum_{j=1}^{d+1}\eta_j\right)
  \prod_{j=1}^{d+1}\dd\eta_j.
\end{aligned}
\eqnum \label{eq:modified-multiplier-scaling-identity}
\]
The inverse dilation therefore replaces the weighted linear constraint with
the ordinary zero-sum constraint and replaces
\(\mathcal D_\omega^{(d-1)}\) by
\(\mathcal R_\omega^{(d-1)}\).  Since
\[
  \mu_v\left(v^{-(d-1)}s\right)
  =\widehat K\left(\frac{v^d}s\right),
\]
the one-variable multiplier in \eqref{eq:modified-multiplier-scaling-identity} is
\[
  \boxed{\qquad
  \widetilde\mu_v(s)
  \eqd\widehat K\left(\frac{v^d}s\right),
  \quad\text{independent of \(\omega\).}
  \qquad}
\]
All \(\omega\)-dependence has been moved into the projective geometry
\(\mathcal R_\omega^{(d-1)}\); the constraint
\(\delta_0(\sum_j\eta_j)\) is independent of \(\omega\).  Each family
\(t\mapsto G_{j,t}\) is defined independently of \(\omega\), and
\(\omega\) merely selects the member \(G_{j,\omega_jv}\).

There is an analogous physical-space description.  Put
\[
  \mathcal P_\omega^{(d-1)}(y_1,\ldots,y_{d+1})
  \eqd
  \det\begin{pmatrix}
    \omega_1y_1&\cdots&\omega_{d+1}y_{d+1}\\
    \omega_1&\cdots&\omega_{d+1}\\
    1&\cdots&1
  \end{pmatrix}.
\]
For \(r>0\) and away from the measure-zero set on which some
\(\omega_j=0\), the changes of variables \(x_j=r\omega_jy_j\) give
\[
\begin{aligned}
  \widetilde{\mathcal S}_{\omega,r}^{(d-1)}
  (H_1,\ldots,H_{d+1})
  ={}&r^{d(d-1)}
  \prod_{j=1}^{d+1}|\omega_j|^{d-1}
  \int_{(\R^{d-1})^{d+1}}\prod_{j=1}^{d+1}H_j(y_j)\\
  &\quad\times
  K\left(
    r^d\mathcal P_\omega^{(d-1)}
    (y_1,\ldots,y_{d+1})
  \right)\\
  &\quad\times
  \delta_0\left(\sum_{j=1}^{d+1}\omega_jy_j\right)
  \prod_{j=1}^{d+1}\dd y_j.
\end{aligned}
\]
Here the inverse dilation replaces the zero-sum constraint by the
\(\omega_j\)-weighted constraint and replaces the determinant by
\(\mathcal P_\omega^{(d-1)}\).  Again, the operator itself receives only
the undilated fibers \(G_{j,r\omega_j}\).

\begin{corollary}[Uniform bounds for the frozen modified operators]
\label{cor:frozen-modified-operator-bounds}
Let \(p_1,\ldots,p_{d+1}\in(1,\infty]\), and put
\[
  \alpha\eqd\sum_{j=1}^{d+1}\frac1{p_j}.
\]
Denote by \(\|T\|_{\vec p}\) the norm of a \((d+1)\)-linear form from
\(\prod_jL^{p_j}(\R^{d-1})\) to \(\C\).  For \(v\ne0\), \(r>0\), and
away from the measure-zero set on which some \(\omega_j=0\), one has
\[
\begin{aligned}
  \bigl\|\widetilde{\mathcal A}_{\omega,v}^{(d-1)}\bigr\|_{\vec p}
  &={}|v|^{(d-1)\alpha}
  \prod_{j=1}^{d+1}|\omega_j|^{(d-1)/p_j}
  \bigl\|\mathcal A_{\omega,\mu_v}^{(d-1)}\bigr\|_{\vec p},\\
  \bigl\|\widetilde{\mathcal S}_{\omega,r}^{(d-1)}\bigr\|_{\vec p}
  &={}r^{(d-1)\alpha}
  \prod_{j=1}^{d+1}|\omega_j|^{(d-1)/p_j}
  \bigl\|\mathcal S_{\omega,K_r}^{(d-1)}\bigr\|_{\vec p}.
\end{aligned}
\]
In particular, uniform bounds for either family of frozen original
operators are equivalent to uniform bounds for the corresponding
scale-normalized modified operators.  Moreover, every frozen modified
operator has the ordinary diagonal translation symmetry, for every
\(z\in\R^{d-1}\):
\[
  \widetilde{\mathcal A}_{\omega,v}^{(d-1)}
  (\Tr_zH_1,\ldots,\Tr_zH_{d+1})
  =
  \widetilde{\mathcal A}_{\omega,v}^{(d-1)}
  (H_1,\ldots,H_{d+1}),
\]
and
\[
  \widetilde{\mathcal S}_{\omega,r}^{(d-1)}
  (\Tr_zH_1,\ldots,\Tr_zH_{d+1})
  =
  \widetilde{\mathcal S}_{\omega,r}^{(d-1)}
  (H_1,\ldots,H_{d+1}).
\]
\end{corollary}

\begin{proof}
The norm identities follow from
\[
  \|H(\cdot/t)\|_{L^p(\R^{d-1})}
  =|t|^{(d-1)/p}\|H\|_{L^p(\R^{d-1})}
\]
and the invertibility of the input dilations.  The multiplier translation
identity follows from the constraint \(\sum_j\eta_j=0\) in
\eqref{eq:modified-multiplier-scaling-identity}.  The space-side identity
follows from \(\sum_j\omega_j=0\): simultaneous translation preserves both
\(\sum_j\omega_jy_j\) and \(\mathcal P_\omega^{(d-1)}(y)\).
\end{proof}

\begin{corollary}[Integrability of the modified operators in the angular and radial parameters]
\label{cor:modified-operator-integrability}
Let \(p_1,\ldots,p_{d+1}\in(1,\infty]\) and \(q\in[1,\infty]\) satisfy
\[
  \frac1q+\sum_{j=1}^{d+1}\frac1{p_j}=1,
  \qquad
  \alpha\eqd\sum_{j=1}^{d+1}\frac1{p_j}.
\]
Suppose that
\[
\begin{aligned}
  \bigl|\widetilde{\mathcal A}_{\omega,v}^{(d-1)}(H_1,\ldots,H_{d+1})\bigr|
  &\le B_{\mathrm{mult}}(v,\omega)\prod_{j=1}^{d+1}\|H_j\|_{p_j},\\
  \bigl|\widetilde{\mathcal S}_{\omega,r}^{(d-1)}(H_1,\ldots,H_{d+1})\bigr|
  &\le B_{\mathrm{sp}}(r,\omega)\prod_{j=1}^{d+1}\|H_j\|_{p_j}.
\end{aligned}
\]
Set
\[
\begin{aligned}
  \widetilde{\mathfrak C}_{\vec p,q}^{\mathrm{mult}}
  &\eqd
  \int_{\mathbb S^{d-1}\cap\vec1^\perp}
  \bigl\||v|^{-(d-1)\alpha}B_{\mathrm{mult}}(v,\omega)
  \bigr\|_{L^q_v(\R)}
  \prod_{j=1}^{d+1}|\omega_j|^{-d/p_j}
  \dd\sigma_\perp(\omega),\\
  \widetilde{\mathfrak C}_{\vec p,q}^{\mathrm{sp}}
  &\eqd
  \int_{\mathbb S^{d-1}\cap\vec1^\perp}
  \bigl\|r^{(d-1)(1-\alpha)}B_{\mathrm{sp}}(r,\omega)
  \bigr\|_{L^q_r(0,\infty)}
  \prod_{j=1}^{d+1}|\omega_j|^{-d/p_j}
  \dd\sigma_\perp(\omega).
\end{aligned}
\]
Then
\[
\begin{aligned}
  \bigl|\mc M_K^{(d)}(g_1,\ldots,g_{d+1})\bigr|
  &\le \frac12\widetilde{\mathfrak C}_{\vec p,q}^{\mathrm{mult}}
  \prod_{j=1}^{d+1}\|g_j\|_{p_j},\\
  \bigl|\Lambda_K^{(d)}(g_1,\ldots,g_{d+1})\bigr|
  &\le \widetilde{\mathfrak C}_{\vec p,q}^{\mathrm{sp}}
  \prod_{j=1}^{d+1}\|g_j\|_{p_j}.
\end{aligned}
\]
In particular, if the bounds in Corollary~\ref{cor:frozen-modified-operator-bounds}
come from coefficients \(C_{\mathrm{mult}}(v,\omega)\) and
\(C_{\mathrm{sp}}(r,\omega)\) for the original frozen operators, the two
integrability conditions reduce respectively to
\[
\begin{aligned}
  &\int_{\mathbb S^{d-1}\cap\vec1^\perp}
  \|C_{\mathrm{mult}}(\cdot,\omega)\|_{L^q(\R)}
  \prod_{j=1}^{d+1}|\omega_j|^{-1/p_j}\dd\sigma_\perp(\omega)<\infty,\\
  &\int_{\mathbb S^{d-1}\cap\vec1^\perp}
  \|r^{d-1}C_{\mathrm{sp}}(r,\omega)\|_{L^q_r(0,\infty)}
  \prod_{j=1}^{d+1}|\omega_j|^{-1/p_j}\dd\sigma_\perp(\omega)<\infty.
\end{aligned}
\]
\end{corollary}

\begin{proof}
For every nonzero \(\omega_j\), the substitution \(t=\omega_jv\) and
\eqref{eq:projective-weighted-norm} give
\[
  \left\||v|^{(d-1)/p_j}
  \|G_{j,\omega_jv}\|_{p_j}\right\|_{L^{p_j}_v(\R)}
  =|\omega_j|^{-d/p_j}\|g_j\|_{p_j}.
\]
The same substitution on \((0,\infty)\) gives
\[
  \left\|r^{(d-1)/p_j}
  \|G_{j,r\omega_j}\|_{p_j}\right\|_{L^{p_j}_r(0,\infty)}
  \le |\omega_j|^{-d/p_j}\|g_j\|_{p_j}.
\]
Apply H\"older's inequality in \(v\), respectively in \(r\), to the two
normalized slicing identities, and then integrate in \(\omega\).  The
last assertion follows from the norm identities in
Corollary~\ref{cor:frozen-modified-operator-bounds}.
\end{proof}

\paragraph{Comparison of the original and modified operators.}
We spell out the four expressions.  The original multiplier-side operator is
\[
\begin{aligned}
  \mathcal A_{\omega,\mu_v}^{(d-1)}
  (H_1,\ldots,H_{d+1})
  ={}&\int_{(\R^{d-1})^{d+1}}
  \prod_{j=1}^{d+1}\widehat H_j(\xi_j)\,
  \mu_v\left(
    \mathcal D_\omega^{(d-1)}
    (\xi_1,\ldots,\xi_{d+1})
  \right)\\
  &\quad\times
  \delta_0\left(\sum_{j=1}^{d+1}\omega_j\xi_j\right)
  \prod_{j=1}^{d+1}\dd\xi_j,
\end{aligned}
\eqnum \label{eq:original-affine-multiplier-form}
\]
whereas the modified operator acting on the fibers is
\[
\begin{aligned}
  \widetilde{\mathcal A}_{\omega,v}^{(d-1)}
  (H_1,\ldots,H_{d+1})
  ={}&|v|^{d-1}
  \int_{(\R^{d-1})^{d+1}}
  \prod_{j=1}^{d+1}\widehat H_j(\eta_j)\,
  \widetilde\mu_v\left(
    \mathcal R_\omega^{(d-1)}
    (\eta_1,\ldots,\eta_{d+1})
  \right)\\
  &\quad\times
  \delta_0\left(\sum_{j=1}^{d+1}\eta_j\right)
  \prod_{j=1}^{d+1}\dd\eta_j.
\end{aligned}
\eqnum \label{eq:modified-affine-multiplier-form}
\]
Thus
\[
\begin{array}{c|c}
  \mathcal A_{\omega,\mu_v}^{(d-1)}
  &\widetilde{\mathcal A}_{\omega,v}^{(d-1)}\\ \hline
  \mathcal D_\omega^{(d-1)}(\xi)
  &\mathcal R_\omega^{(d-1)}(\eta)\\
  \sum_j\omega_j\xi_j=0
  &\sum_j\eta_j=0\\
  \mu_v(s)=\widehat K(v/s)
  &\widetilde\mu_v(s)=\widehat K(v^d/s)\\
  1&|v|^{d-1}
\end{array}
\]
in the determinant, constraint, one-variable multiplier, and exterior
Jacobian, respectively.  The two determinants are related by
\[
  \mathcal R_\omega^{(d-1)}(\eta_1,\ldots,\eta_{d+1})
  =\mathcal D_\omega^{(d-1)}
  (\eta_1/\omega_1,\ldots,\eta_{d+1}/\omega_{d+1}),
\]
This identity is understood away from the measure-zero set on which some
\(\omega_j=0\).  The essential gain of this normalization is that its
constraint is the ordinary zero-sum constraint, uniformly in \(\omega\).

On the space side, the original expression is
\[
\begin{aligned}
  \mathcal S_{\omega,K_r}^{(d-1)}
  (H_1,\ldots,H_{d+1})
  ={}&\int_{(\R^{d-1})^{d+1}}
  \prod_{j=1}^{d+1}H_j(x_j)\,
  K\left(
    r\mathcal D_\omega^{(d-1)}
    (x_1,\ldots,x_{d+1})
  \right)\\
  &\quad\times
  \delta_0\left(\sum_{j=1}^{d+1}x_j\right)
  \prod_{j=1}^{d+1}\dd x_j,
\end{aligned}
\eqnum \label{eq:original-affine-space-form}
\]
while the modified fiber operator is
\[
\begin{aligned}
  \widetilde{\mathcal S}_{\omega,r}^{(d-1)}
  (H_1,\ldots,H_{d+1})
  ={}&r^{d(d-1)}
  \prod_{j=1}^{d+1}|\omega_j|^{d-1}
  \int_{(\R^{d-1})^{d+1}}\prod_{j=1}^{d+1}H_j(y_j)\\
  &\quad\times
  K\left(
    r^d\mathcal P_\omega^{(d-1)}
    (y_1,\ldots,y_{d+1})
  \right)
  \delta_0\left(\sum_{j=1}^{d+1}\omega_jy_j\right)
  \prod_{j=1}^{d+1}\dd y_j.
\end{aligned}
\eqnum \label{eq:modified-affine-space-form}
\]
Here the zero-sum plane becomes the \(\omega\)-weighted plane, the
determinant becomes \(\mathcal P_\omega^{(d-1)}\), the kernel scale changes
from \(r\) to \(r^d\), and the change of variables contributes the
displayed Jacobian.

The original and modified operators are therefore \emph{not equal as
multilinear forms evaluated on the same functions}.  Their exact relation
is instead
\[
\begin{aligned}
  \widetilde{\mathcal A}_{\omega,v}^{(d-1)}(H_1,\ldots,H_{d+1})
  &=
  \mathcal A_{\omega,\mu_v}^{(d-1)}
  \bigl(H_1(\cdot/(\omega_1v)),\ldots,
        H_{d+1}(\cdot/(\omega_{d+1}v))\bigr),\\
  \widetilde{\mathcal S}_{\omega,r}^{(d-1)}(H_1,\ldots,H_{d+1})
  &=
  \mathcal S_{\omega,K_r}^{(d-1)}
  \bigl(H_1(\cdot/(r\omega_1)),\ldots,
        H_{d+1}(\cdot/(r\omega_{d+1}))\bigr).
\end{aligned}
\]
For fixed \(v\ne0\), \(r>0\), and \(\omega_j\ne0\), these input dilations
are invertible.  Hence the original and modified operators contain the
same analytic information, with the corresponding \(L^{p_j}\)-scaling
factors, but only the modified operators act directly on the fibers.

\subsection{Concentration in the last variable.}

We now test both models on functions concentrated near prescribed last
coordinates.  Let \(B_\varepsilon(\zeta)=(\zeta-\varepsilon,
\zeta+\varepsilon)\), let \(h_j\in\Sch(\R^{d-1})\), and put
\[
  g_j^{(\varepsilon)}(x,t)
  \eqd
  \varepsilon^{-1/p_j}h_j(x)
  \1_{B_\varepsilon(\zeta_j)}(t).
  \eqnum \label{eq:concentrated-test-functions}
\]
Then
\[
  \|g_j^{(\varepsilon)}\|_{L^{p_j}(\R^d)}
  =2^{1/p_j}\|h_j\|_{L^{p_j}(\R^{d-1})},
\]
so these norms are independent of \(\varepsilon\).  The harmless factor
\(2^{1/p_j}\) disappears if \((2\varepsilon)^{-1/p_j}\) is used in
\eqref{eq:concentrated-test-functions}.  Write
\[
  \alpha\eqd\sum_{j=1}^{d+1}\frac1{p_j},
  \qquad
  \zeta=(\zeta_1,\ldots,\zeta_{d+1}).
\]

On the space side, the zero-sum constraint forces
\(\sum_jt_j=0\).  Consequently, if \(\sum_j\zeta_j\ne0\), then
\[
  \Lambda_K^{(d)}
  (g_1^{(\varepsilon)},\ldots,g_{d+1}^{(\varepsilon)})=0
\]
for all sufficiently small \(\varepsilon\).  Suppose instead that
\(\zeta\in\vec1^\perp\).  The symmetric slicing formula gives the exact
identity
\[
\begin{aligned}
  &\Lambda_K^{(d)}
  (g_1^{(\varepsilon)},\ldots,g_{d+1}^{(\varepsilon)})\\
  &\qquad=
  \varepsilon^{-\alpha}
  \int_{\vec1^\perp}
  \prod_{j=1}^{d+1}
  \1_{B_\varepsilon(\zeta_j)}(t_j)\,
  \mathcal S_{t,K}^{(d-1)}(h_1,\ldots,h_{d+1})
  \dd\mc L^d(t)\\
  &\qquad=
  \varepsilon^{d-\alpha}
  \int_{\vec1^\perp}
  \prod_{j=1}^{d+1}\1_{(-1,1)}(s_j)\,
  \mathcal S_{\zeta+\varepsilon s,K}^{(d-1)}
  (h_1,\ldots,h_{d+1})
  \dd\mc L^d(s).
\end{aligned}
\eqnum \label{eq:space-concentration-identity}
\]
The power \(\varepsilon^d\) is the volume of the intersection of the
\((d+1)\)-dimensional box with the \(d\)-dimensional constraint plane.
For a regularized kernel, continuity of the sliced form therefore gives
\[
  \Lambda_K^{(d)}
  (g_1^{(\varepsilon)},\ldots,g_{d+1}^{(\varepsilon)})
  =
  c_d\varepsilon^{d-\alpha}
  \mathcal S_{\zeta,K}^{(d-1)}(h_1,\ldots,h_{d+1})
  +o(\varepsilon^{d-\alpha}),
  \eqnum \label{eq:space-concentration-asymptotic}
\]
where
\[
  c_d\eqd
  \int_{\vec1^\perp}
  \prod_{j=1}^{d+1}\1_{(-1,1)}(s_j)\dd\mc L^d(s).
\]
As usual, the statement for a distributional \(K\) is read after
regularization.

The corresponding normalized fiber is
\[
  G_{j,t}^{(\varepsilon)}(x)
  =
  \varepsilon^{-1/p_j}
  h_j(tx)\1_{B_\varepsilon(\zeta_j)}(t).
\]
If \(t_j=r\omega_j\ne0\) for every \(j\), then the modified space-side
operator acts on these fibers without any further change of input:
\[
  \widetilde{\mathcal S}_{\omega,r}^{(d-1)}
  \bigl(
    G_{1,t_1}^{(\varepsilon)},\ldots,
    G_{d+1,t_{d+1}}^{(\varepsilon)}
  \bigr)
  =
  \varepsilon^{-\alpha}
  \prod_{j=1}^{d+1}\1_{B_\varepsilon(\zeta_j)}(t_j)\,
  \mathcal S_{\omega,K_r}^{(d-1)}
  (h_1,\ldots,h_{d+1}).
\]
Indeed,
\(G_{j,t_j}^{(\varepsilon)}(x/t_j)
=\varepsilon^{-1/p_j}h_j(x)
\1_{B_\varepsilon(\zeta_j)}(t_j)\).

On the multiplier side, the same concentration forces
\(\omega_jv\in B_\varepsilon(\zeta_j)\) for every \(j\).  Hence it again
vanishes for all sufficiently small \(\varepsilon\) unless
\(\zeta\in\vec1^\perp\).  Assume that
\(\zeta\in\vec1^\perp\setminus\{0\}\).  Substitution in the spherical
coarea formula gives
\[
\begin{aligned}
  &\mc M_K^{(d)}
  (g_1^{(\varepsilon)},\ldots,g_{d+1}^{(\varepsilon)})\\
  &\quad=
  \frac{\varepsilon^{-\alpha}}2
  \int_\R\int_{\mathbb S^{d-1}\cap\vec1^\perp}
  \prod_{j=1}^{d+1}
  \1_{B_\varepsilon(\zeta_j)}(\omega_jv)\,
  \mathcal A_{\omega,\mu_v}^{(d-1)}
  (h_1,\ldots,h_{d+1})
  \dd\sigma_\perp(\omega)\dd v.
\end{aligned}
\eqnum \label{eq:multiplier-concentration-identity}
\]
The pairs \((v,\omega)\) and \((-v,-\omega)\) describe the same point
\(t=v\omega\in\vec1^\perp\), and
\[
  \mathcal A_{-\omega,\mu_{-v}}^{(d-1)}
  =\mathcal A_{\omega,\mu_v}^{(d-1)}.
\]
Since
\(\dd\mc L^d(t)=|v|^{d-1}\dd v\dd\sigma_\perp(\omega)\), the factor
\(1/2\) removes this double covering and \eqref{eq:multiplier-concentration-identity}
becomes
\[
\begin{aligned}
  &\mc M_K^{(d)}
  (g_1^{(\varepsilon)},\ldots,g_{d+1}^{(\varepsilon)})\\
  &\quad=
  \varepsilon^{d-\alpha}
  \int_{\vec1^\perp}
  \prod_{j=1}^{d+1}\1_{(-1,1)}(s_j)\,
  |\zeta+\varepsilon s|^{-(d-1)}\\
  &\qquad\qquad\times
  \mathcal A_{
    (\zeta+\varepsilon s)/|\zeta+\varepsilon s|,
    \mu_{|\zeta+\varepsilon s|}
  }^{(d-1)}
  (h_1,\ldots,h_{d+1})
  \dd\mc L^d(s).
\end{aligned}
\eqnum \label{eq:multiplier-concentration-polar-identity}
\]
Thus, for a regularized kernel,
\[
\begin{aligned}
  &\mc M_K^{(d)}
  (g_1^{(\varepsilon)},\ldots,g_{d+1}^{(\varepsilon)})\\
  &\qquad=
  c_d\varepsilon^{d-\alpha}|\zeta|^{-(d-1)}
  \mathcal A_{\zeta/|\zeta|,\mu_{|\zeta|}}^{(d-1)}
  (h_1,\ldots,h_{d+1})
  +o(\varepsilon^{d-\alpha}).
\end{aligned}
\eqnum \label{eq:multiplier-concentration-asymptotic}
\]
The additional factor \(|\zeta|^{-(d-1)}\), absent from
\eqref{eq:space-concentration-asymptotic}, is precisely the polar-coordinate
Jacobian in the multiplier-side coarea formula.

In terms of the modified multiplier operator, if
\(t_j=\omega_jv\ne0\) for every \(j\), then
\[
  \widetilde{\mathcal A}_{\omega,v}^{(d-1)}
  \bigl(
    G_{1,t_1}^{(\varepsilon)},\ldots,
    G_{d+1,t_{d+1}}^{(\varepsilon)}
  \bigr)
  =
  \varepsilon^{-\alpha}
  \prod_{j=1}^{d+1}\1_{B_\varepsilon(\zeta_j)}(t_j)\,
  \mathcal A_{\omega,\mu_v}^{(d-1)}
  (h_1,\ldots,h_{d+1}).
\]
Thus in both models the modified operator acts directly on the concentrated
fibers; its built-in inverse dilation cancels the \(t_jx\) appearing in
\(G_{j,t_j}^{(\varepsilon)}\).
The excluded parameter set on which some
\(t_j=0\) has measure zero and may be handled by a limiting convention.

In the scale-invariant case
\(\alpha=1\), both forms are of size \(\varepsilon^{d-1}\) for this test.
Maintaining the individual \(L^{p_j}\)-norms therefore does not produce a
nonzero limiting form when \(d\ge2\); to extract the indicated slice one
must additionally divide the form by \(\varepsilon^{d-1}\).

\paragraph{Scale-invariant diagonal packet series.}
The preceding construction can be repeated at many separated scales.
For every sequence index \(n\), choose
\[
  \zeta^{(n)}
  =(\zeta_{1,n},\ldots,\zeta_{d+1,n})\in\vec1^\perp
  \qquad\text{and}\qquad r_n>0.
\]
We impose the no-mixed-resonance condition
\[
  \left|\sum_{j=1}^{d+1}\zeta_{j,n_j}\right|
  >
  \sum_{j=1}^{d+1}r_{n_j}
  \quad\text{whenever the indices \(n_1,\ldots,n_{d+1}\)
  are not all equal.}
  \eqnum \label{eq:no-mixed-resonance-condition}
\]
It follows that
\[
  \vec1^\perp\cap
  \prod_{j=1}^{d+1}B_{r_{n_j}}(\zeta_{j,n_j})
  =\varnothing
\]
unless \(n_1=\cdots=n_{d+1}\).  Thus every mixed term in the expansion of
the full form vanishes exactly.

There is a natural scale-invariant way to arrange \((\mathsf D)\).  Choose
\(\zeta^\circ\in\vec1^\perp\) such that no proper nonempty subsum of its
coordinates vanishes, take a sufficiently lacunary sequence
\(\lambda_n=L^n\), and set
\[
  \zeta^{(n)}=\lambda_n\zeta^\circ,
  \qquad
  r_n=\delta\lambda_n.
  \eqnum \label{eq:lacunary-scale-choice}
\]
If \(L\) is sufficiently large and then \(\delta>0\) is sufficiently
small, both \((\mathsf D)\) and the pairwise disjointness of the intervals
\(B_{r_n}(\zeta_{j,n})\), for each fixed \(j\), hold.  Indeed, in any
off-diagonal sum the terms at the largest scale contribute a proper
subsum of \(\zeta^\circ\); lacunarity makes all lower scales too small to
cancel it.

First prescribe the packets in the original functions:
\[
  g_j(x,t)
  \eqd
  \sum_n r_n^{-1/p_j}h_{j,n}(x)
  \1_{B_{r_n}(\zeta_{j,n})}(t).
  \eqnum \label{eq:original-packet-series}
\]
Disjointness gives
\[
  \|g_j\|_{L^{p_j}(\R^d)}^{p_j}
  =2\sum_n\|h_{j,n}\|_{L^{p_j}(\R^{d-1})}^{p_j}.
  \eqnum \label{eq:original-packet-series-norms}
\]
For \(p_j=\infty\), the right-hand side is replaced by the corresponding
supremum.  Define the diagonal block operator
\[
\begin{aligned}
  \mathfrak B_n(h_{1,n},\ldots,h_{d+1,n})
  \eqd{}&
  \frac{r_n^{-\alpha}}2
  \int_\R\int_{\mathbb S^{d-1}\cap\vec1^\perp}
  \prod_{j=1}^{d+1}
  \1_{B_{r_n}(\zeta_{j,n})}(\omega_jv)\\
  &\qquad\times
  \mathcal A_{\omega,\mu_v}^{(d-1)}
  (h_{1,n},\ldots,h_{d+1,n})
  \dd\sigma_\perp(\omega)\dd v.
\end{aligned}
\eqnum \label{eq:diagonal-block-operator}
\]
Condition \((\mathsf D)\) gives the exact diagonalization
\[
  \mc M_K^{(d)}(g_1,\ldots,g_{d+1})
  =\sum_n
  \mathfrak B_n(h_{1,n},\ldots,h_{d+1,n}).
  \eqnum \label{eq:diagonalization-of-full-form}
\]
Consequently, an estimate
\[
  |\mc M_K^{(d)}(g_1,\ldots,g_{d+1})|
  \le C_{\vec p}\prod_{j=1}^{d+1}
  \|g_j\|_{L^{p_j}(\R^d)}
  \eqnum \label{eq:full-form-bound}
\]
implies the diagonal vector-valued inequality
\[
  \left|
  \sum_n\mathfrak B_n
  (h_{1,n},\ldots,h_{d+1,n})
  \right|
  \le
  2^\alpha C_{\vec p}
  \prod_{j=1}^{d+1}
  \left(
    \sum_n
    \|h_{j,n}\|_{L^{p_j}(\R^{d-1})}^{p_j}
  \right)^{1/p_j}.
  \eqnum \label{eq:diagonal-vector-valued-bound}
\]
This is the basic family of translation-invariant, \(A\)-type bounds
forced by the full estimate.

The geometry and scale of each block are transparent after setting
\(t=\zeta^{(n)}+r_ns\):
\[
\begin{aligned}[t]
  \mathfrak B_n(h_{1,n},\ldots,h_{d+1,n})
  =&r_n^{d-\alpha}
  \int_{\vec1^\perp}
  \prod_{j=1}^{d+1}\1_{(-1,1)}(s_j)\,
  |\zeta^{(n)}+r_ns|^{-(d-1)}\\
  &\quad\times
  \mathcal A_{
    (\zeta^{(n)}+r_ns)/|\zeta^{(n)}+r_ns|,
    \mu_{|\zeta^{(n)}+r_ns|}
  }^{(d-1)}
  (h_{1,n},\ldots,h_{d+1,n})
  \dd\mc L^d(s).
\end{aligned}
\eqnum \label{eq:diagonal-block-polar-form}
\]
Under the scale-invariant choice \((\mathsf L)\), this becomes
\[
  \mathfrak B_n
  =\delta^{d-\alpha}\lambda_n^{1-\alpha}
  \mathfrak T_{\delta,\lambda_n},
  \eqnum \label{eq:diagonal-block-scaling}
\]
where
\[
\begin{aligned}
  \mathfrak T_{\delta,\lambda}(h_1,\ldots,h_{d+1})
  \eqd{}&
  \int_{\vec1^\perp}
  \prod_{j=1}^{d+1}\1_{(-1,1)}(s_j)\,
  |\zeta^\circ+\delta s|^{-(d-1)}\\
  &\quad\times
  \mathcal A_{
    (\zeta^\circ+\delta s)/|\zeta^\circ+\delta s|,
    \mu_{\lambda|\zeta^\circ+\delta s|}
  }^{(d-1)}
  (h_1,\ldots,h_{d+1})
  \dd\mc L^d(s).
\end{aligned}
\]
In the critical case \(\alpha=1\), the scale factor
\(\lambda_n^{1-\alpha}\) disappears.  Thus \eqref{eq:diagonal-vector-valued-bound} implies
\[
  \left|
  \sum_n\mathfrak T_{\delta,\lambda_n}
  (h_{1,n},\ldots,h_{d+1,n})
  \right|
  \lesssim_{\vec p}
  C_{\vec p}\delta^{-(d-1)}
  \prod_{j=1}^{d+1}
  \left(
    \sum_n\|h_{j,n}\|_{p_j}^{p_j}
  \right)^{1/p_j}.
  \eqnum \label{eq:translation-vector-valued-bound}
\]
Moreover, for a regularized kernel,
\[
  \mathfrak T_{\delta,\lambda}
  \longrightarrow
  c_d|\zeta^\circ|^{-(d-1)}
  \mathcal A_{
    \zeta^\circ/|\zeta^\circ|,
    \mu_{\lambda|\zeta^\circ|}
  }^{(d-1)}
  \quad\text{as }\delta\to0.
  \eqnum \label{eq:block-operator-pointwise-limit}
\]

One may instead prescribe the series directly in the normalized fibers.
For an interval not meeting zero, put
\[
  I_{j,n}
  \eqd
  \int_{B_{r_n}(\zeta_{j,n})}
  |t|^{d-1}\dd t,
  \qquad
  a_{j,n}\eqd I_{j,n}^{-1/p_j},
\]
and define
\[
  G_{j,t}(x)
  \eqd
  \sum_n a_{j,n}H_{j,n}(x)
  \1_{B_{r_n}(\zeta_{j,n})}(t).
  \eqnum \label{eq:normalized-fiber-series}
\]
The weighted fiber identity \eqref{eq:projective-weighted-norm} gives exactly
\[
  \|g_j\|_{L^{p_j}(\R^d)}^{p_j}
  =
  \sum_n\|H_{j,n}\|_{L^{p_j}(\R^{d-1})}^{p_j},
  \qquad
  g_j(x,t)=G_{j,t}(x/t).
  \eqnum \label{eq:normalized-fiber-series-norms}
\]
In this parametrization define
\[
\begin{aligned}
  \widetilde{\mathfrak B}_n
  (H_{1,n},\ldots,H_{d+1,n})
  \eqd{}&
  \frac12\left(\prod_{j=1}^{d+1}a_{j,n}\right)
  \int_\R\int_{\mathbb S^{d-1}\cap\vec1^\perp}
  \prod_{j=1}^{d+1}
  \1_{B_{r_n}(\zeta_{j,n})}(\omega_jv)\\
  &\qquad\times
  \widetilde{\mathcal A}_{\omega,v}^{(d-1)}
  (H_{1,n},\ldots,H_{d+1,n})
  \dd\sigma_\perp(\omega)\dd v.
\end{aligned}
\eqnum \label{eq:normalized-diagonal-block-operator}
\]
Now the full form again diagonalizes exactly:
\[
  \mc M_K^{(d)}(g_1,\ldots,g_{d+1})
  =
  \sum_n\widetilde{\mathfrak B}_n
  (H_{1,n},\ldots,H_{d+1,n}),
\]
and \((\mathsf M_{\vec p})\) implies
\[
  \left|
  \sum_n\widetilde{\mathfrak B}_n
  (H_{1,n},\ldots,H_{d+1,n})
  \right|
  \le C_{\vec p}
  \prod_{j=1}^{d+1}
  \left(
    \sum_n\|H_{j,n}\|_{p_j}^{p_j}
  \right)^{1/p_j}.
  \eqnum \label{eq:normalized-diagonal-vector-valued-bound}
\]
This is the cleanest formulation if one wants operators acting only on
fibers.  By Corollary~\ref{cor:frozen-modified-operator-bounds}, every
constituent, and hence every block average, has the same ordinary diagonal
translation invariance.

The conclusion has an important limitation.  The full bound
\eqref{eq:full-form-bound} gives the uniform diagonal estimate
\eqref{eq:diagonal-vector-valued-bound}, or equivalently
\eqref{eq:normalized-diagonal-vector-valued-bound}, for the block-averaged operators.  It does
not, by this argument alone, give a uniform pointwise estimate for
\(\mathcal A_{\omega,\mu_v}^{(d-1)}\).  Indeed, in the critical case the
block has transverse volume \(\delta^{d-1}\), and normalizing it to converge
to the pointwise operator produces the loss
\(\delta^{-(d-1)}\) in \eqref{eq:translation-vector-valued-bound}.  Adding more diagonal scales does
not remove this loss: when \(\alpha=1\), both sides of the diagonal
inequality grow linearly under repetition.  What the full estimate forces
without further input is therefore an \(\ell^{p_j}\)-valued bound for
scale-separated averages of translation-invariant \(A\)-type operators,
not a \(\delta\)-uniform bound for an individual slice.

\begin{corollary}[The multiplier form in dimension two]
\label{cor:multiplier-affine-trick-plane}
Let $d=2$, let $K(t)=\pv(1/t)$, and let
$p_1,p_2,p_3\in(2,\infty)$ satisfy
\(1/p_1+1/p_2+1/p_3=1\).  Then
\[
  \bigl|\mc M_K^{(2)}(g_1,g_2,g_3)\bigr|
  \lesssim_{\vec p}
  \prod_{j=1}^3\|g_j\|_{L^{p_j}(\R^2)}.
\]
\end{corollary}

\begin{proof}
The uniform bilinear Hilbert transform estimate
of Theorem~\ref{thm:uniform-bht}, after the affine changes of frequency
variables determined by $\omega$, gives
\[
  \bigl|\mathcal A_{\omega,\sgn}^{(1)}(h_1,h_2,h_3)\bigr|
  \lesssim_{\vec p}
  \prod_{j=1}^3\|h_j\|_{L^{p_j}(\R)}
\]
uniformly in $\omega\in\mathbb S^1\cap\vec1^\perp$.  Since
$\mu_v=-i\pi\sgn(v)\sgn$, this is the uniform bound required by
Corollary~\ref{cor:uniform-multiplier-sliced-bound} with $q=\infty$.
Finally,
\[
  \int_{\mathbb S^1\cap\vec1^\perp}
  \prod_{j=1}^3|\omega_j|^{-1/p_j}\,
  \dd\sigma_\perp(\omega)<\infty
\]
because $p_j>2$.  The conclusion follows.
\end{proof}
